\documentclass[10pt,a4paper]{amsart}
\usepackage{amsmath,amssymb,amsthm,graphicx}
% For commutative diagrams we use Xy-pic
\usepackage[all]{xy}

% We use 2cell for 2-commutative diagrams.
\xyoption{2cell}
% Theorem environments.
%
\theoremstyle{plain}
\newtheorem{theorem}[subsection]{Theorem}
\newtheorem{proposition}[subsection]{Proposition}
\newtheorem{lemma}[subsection]{Lemma}

\theoremstyle{definition}
\newtheorem{definition}[subsection]{Definition}
\newtheorem{example}[subsection]{Example}
\newtheorem{exercise}[subsection]{Exercise}
\newtheorem{situation}[subsection]{Situation}

\theoremstyle{remark}
\newtheorem{remark}[subsection]{Remark}
\newtheorem{remarks}[subsection]{Remarks}

\numberwithin{equation}{subsection}

% Macros
%
\def\lim{\mathop{\mathrm{lim}}\nolimits}
\def\colim{\mathop{\mathrm{colim}}\nolimits}
\def\Spec{\mathop{\mathrm{Spec}}}
\def\Hom{\mathop{\mathrm{Hom}}\nolimits}
\def\Ext{\mathop{\mathrm{Ext}}\nolimits}
\def\SheafHom{\mathop{\mathcal{H}\!\mathit{om}}\nolimits}
\def\SheafExt{\mathop{\mathcal{E}\!\mathit{xt}}\nolimits}
\def\Sch{\mathit{Sch}}
\def\Mor{\mathop{\mathrm{Mor}}\nolimits}
\def\Ob{\mathop{\mathrm{Ob}}\nolimits}
\def\Sh{\mathop{\mathit{Sh}}\nolimits}
\def\NL{\mathop{N\!L}\nolimits}
\def\CH{\mathop{\mathrm{CH}}\nolimits}
\def\proetale{{pro\text{-}\acute{e}tale}}
\def\etale{{\acute{e}tale}}
\def\QCoh{\mathit{QCoh}}
\def\Ker{\mathop{\mathrm{Ker}}}
\def\Im{\mathop{\mathrm{Im}}}
\def\Coker{\mathop{\mathrm{Coker}}}
\def\Coim{\mathop{\mathrm{Coim}}}

% Boxtimes
%
\DeclareMathSymbol{\boxtimes}{\mathbin}{AMSa}{"02}

%
% Macros for moduli stacks/spaces
%
\def\QCohstack{\mathcal{QC}\!\mathit{oh}}
\def\Cohstack{\mathcal{C}\!\mathit{oh}}
\def\Spacesstack{\mathcal{S}\!\mathit{paces}}
\def\Quotfunctor{\mathrm{Quot}}
\def\Hilbfunctor{\mathrm{Hilb}}
\def\Curvesstack{\mathcal{C}\!\mathit{urves}}
\def\Polarizedstack{\mathcal{P}\!\mathit{olarized}}
\def\Complexesstack{\mathcal{C}\!\mathit{omplexes}}
% \Pic is the operator that assigns to X its picard group, usage \Pic(X)
% \Picardstack_{X/B} denotes the Picard stack of X over B
% \Picardfunctor_{X/B} denotes the Picard functor of X over B
\def\Pic{\mathop{\mathrm{Pic}}\nolimits}
\def\Picardstack{\mathcal{P}\!\mathit{ic}}
\def\Picardfunctor{\mathrm{Pic}}
\def\Deformationcategory{\mathcal{D}\!\mathit{ef}}
\usepackage{cleveref}
\crefname{originallemma}{original lemma}{original lemmas}
\Crefname{originallemma}{Original lemma}{Original lemmas}
\makeatletter
\newlabel{lemma-compare-cohomology-point}{{03QU}{source}{}{sourcepointer-1}{}}
\newlabel{lemma-compare-cohomology-point@cref}{{[originallemma][0][]03QU}{[1][0][]source}{}{}{}}
\makeatother
\crefname{originaltheorem}{original theorem}{original theorems}
\Crefname{originaltheorem}{Original theorem}{Original theorems}
\makeatletter
\newlabel{theorem-brauer-delta}{{03R7}{source}{}{sourcepointer-2}{}}
\newlabel{theorem-brauer-delta@cref}{{[originaltheorem][0][]03R7}{[1][0][]source}{}{}{}}
\makeatother
\crefname{originalremark}{original remark}{original remarks}
\Crefname{originalremark}{Original remark}{Original remarks}
\makeatletter
\newlabel{example-sheaves-point}{{03QV}{source}{}{sourcepointer-3}{}}
\newlabel{example-sheaves-point@cref}{{[originalremark][0][]03QV}{[1][0][]source}{}{}{}}
\makeatother
\crefname{originaltheorem}{original theorem}{original theorems}
\Crefname{originaltheorem}{Original theorem}{Original theorems}
\makeatletter
\newlabel{theorem-colimit}{{09YQ}{source}{}{sourcepointer-4}{}}
\newlabel{theorem-colimit@cref}{{[originaltheorem][0][]09YQ}{[1][0][]source}{}{}{}}
\makeatother
\crefname{originallemma}{original lemma}{original lemmas}
\Crefname{originallemma}{Original lemma}{Original lemmas}
\makeatletter
\newlabel{lemma-kummer-sequence}{{03PL}{source}{}{sourcepointer-5}{}}
\newlabel{lemma-kummer-sequence@cref}{{[originallemma][0][]03PL}{[1][0][]source}{}{}{}}
\makeatother
\crefname{originaltheorem}{original theorem}{original theorems}
\Crefname{originaltheorem}{Original theorem}{Original theorems}
\makeatletter
\newlabel{theorem-picard-group}{{03P8}{source}{}{sourcepointer-6}{}}
\newlabel{theorem-picard-group@cref}{{[originaltheorem][0][]03P8}{[1][0][]source}{}{}{}}
\makeatother
\crefname{originallemma}{original lemma}{original lemmas}
\Crefname{originallemma}{Original lemma}{Original lemmas}
\makeatletter
\newlabel{lemma-reduce-to-l-group-higher}{{0DV9}{source}{}{sourcepointer-7}{}}
\newlabel{lemma-reduce-to-l-group-higher@cref}{{[originallemma][0][]0DV9}{[1][0][]source}{}{}{}}
\makeatother
\crefname{originalsection}{original section}{original sections}
\Crefname{originalsection}{Original section}{Original sections}
\makeatletter
\newlabel{section-artin-schreier}{{0A3J}{source}{}{sourcepointer-8}{}}
\newlabel{section-artin-schreier@cref}{{[originalsection][0][]0A3J}{[1][0][]source}{}{}{}}
\makeatother
\crefname{originalremark}{original remark}{original remarks}
\Crefname{originalremark}{Original remark}{Original remarks}
\makeatletter
\newlabel{remark-special-case-fpqc-cohomology-quasi-coherent}{{03P6}{source}{}{sourcepointer-9}{}}
\newlabel{remark-special-case-fpqc-cohomology-quasi-coherent@cref}{{[originalremark][0][]03P6}{[1][0][]source}{}{}{}}
\makeatother
\crefname{originallemma}{original lemma}{original lemmas}
\Crefname{originallemma}{Original lemma}{Original lemmas}
\makeatletter
\newlabel{coherent-lemma-quasi-coherent-affine-cohomology-zero}{{01XB}{source}{}{sourcepointer-10}{}}
\newlabel{coherent-lemma-quasi-coherent-affine-cohomology-zero@cref}{{[originallemma][0][]01XB}{[1][0][]source}{}{}{}}
\makeatother
\crefname{originallemma}{original lemma}{original lemmas}
\Crefname{originallemma}{Original lemma}{Original lemmas}
\makeatletter
\newlabel{lemma-profinite-group-cohomology-torsion}{{0DV3}{source}{}{sourcepointer-11}{}}
\newlabel{lemma-profinite-group-cohomology-torsion@cref}{{[originallemma][0][]0DV3}{[1][0][]source}{}{}{}}
\makeatother
\title{Serre's vanishing theorem for Galois cohomology under the vanishing of Brauer groups}
\author{The Stacks Project Authors}
\date{stacks-project:etale-cohomology.tex:local-snapshot-20261011; GNU FDL 1.2 or later}
\begin{document}
\maketitle
\noindent\textit{Editorial note.} All delivered mathematical text is the author's own, copied from the pinned Stacks Project chapter file etale-cohomology.tex; no statement or proof step was rewritten, reordered or removed.\par\medskip
\noindent\textit{Editorial note.} The statement environment of this proposition contains a commented reference block, \textbackslash{}begin\{reference\}\textbackslash{}cite[Chapter II, Section 3, Proposition 5]\{SerreGaloisCohomology\}\textbackslash{}end\{reference\}, which records the attribution of the result. That block is excluded from the counted spans by splitting the statement, so that no bibliography entry has to be supplied for the card; nothing mathematical is dropped by the split, and the attribution is recorded here instead.\par\medskip
\noindent\textit{Editorial note.} The proof is the author's complete argument for both assertions: the translation to etale cohomology, the two characteristic cases and the final dimension argument are all inside the counted span.\par\medskip
\noindent\textit{Editorial note.} The checked chapter list for same-name collisions is the set of stacks chapters staged in this repository: chow, cohomology, divisors, etale-cohomology, flat, more-algebra, more-morphisms, morphisms, spaces-more-morphisms, varieties together with this card's chapter etale-cohomology. Every labelled reference inside the counted spans is carried in assembly.external\_references as a non-load-bearing navigation pointer.\par\medskip
\noindent\textit{Editorial note.} The preamble is the chapter's own native preamble excerpt already used for this source: the Xy-pic setup (native preamble lines 14-18) and the native theorem declarations with the mathematical macros (native preamble lines 154-217).\par\medskip
\noindent\textit{Allowed standard prerequisites (editorial).}\begin{itemize}
\item Comparison between Galois cohomology and etale cohomology of Spec(K). Same chapter, the comparison result identifying H\textasciicircum{}q(Gal(K\textasciicircum{}sep/K), M) with the etale cohomology of Spec(K) with coefficients in the corresponding sheaf. Supplies the translation that lets the whole proposition be proved as a statement about etale cohomology.
\item Kummer sequence and Artin-Schreier sequence. Same chapter, the short exact Kummer sequence for l invertible in K and the Artin-Schreier sequence in characteristic p. Supplies the two exact sequences whose long exact cohomology sequences give the vanishing, according to whether l differs from or equals the characteristic.
\item Vanishing of the Brauer group of finite extensions (the hypothesis of the statement). The hypothesis Br(K') = 0 for every finite extension K' of K, together with the identification of Br(K') with H\textasciicircum{}2 of the Galois group with coefficients in the multiplicative group. Supplies the input that kills the second cohomology of the multiplicative group, hence the Brauer group, in the Kummer case.
\end{itemize}\medskip
\ifdefined\hypertarget\hypertarget{sourcepointer-1}{}\fi
\noindent\textit{Source reference.} Original lemma 03QU: sources/stacks-etale-cohomology/pinned-original-source.tex, source lines 9100--9100. Navigation pointer invoked at proof line 10711; the counted span contains the pointer only and the cited result is not reproduced inside it, so the pointer is registered as a non-load-bearing navigation stub in the established form for Stacks material. No bibliography citation is involved.\par\medskip
\ifdefined\hypertarget\hypertarget{sourcepointer-2}{}\fi
\noindent\textit{Source reference.} Original theorem 03R7: sources/stacks-etale-cohomology/pinned-original-source.tex, source lines 9415--9415. Navigation pointer invoked at proof line 10712; the counted span contains the pointer only and the cited result is not reproduced inside it, so the pointer is registered as a non-load-bearing navigation stub in the established form for Stacks material. No bibliography citation is involved.\par\medskip
\ifdefined\hypertarget\hypertarget{sourcepointer-3}{}\fi
\noindent\textit{Source reference.} Original remark 03QV: sources/stacks-etale-cohomology/pinned-original-source.tex, source lines 9122--9122. Navigation pointer invoked at proof line 10713; the counted span contains the pointer only and the cited result is not reproduced inside it, so the pointer is registered as a non-load-bearing navigation stub in the established form for Stacks material. No bibliography citation is involved.\par\medskip
\ifdefined\hypertarget\hypertarget{sourcepointer-4}{}\fi
\noindent\textit{Source reference.} Original theorem 09YQ: sources/stacks-etale-cohomology/pinned-original-source.tex, source lines 7400--7400. Navigation pointer invoked at proof line 10633; the counted span contains the pointer only and the cited result is not reproduced inside it, so the pointer is registered as a non-load-bearing navigation stub in the established form for Stacks material. No bibliography citation is involved.\par\medskip
\ifdefined\hypertarget\hypertarget{sourcepointer-5}{}\fi
\noindent\textit{Source reference.} Original lemma 03PL: sources/stacks-etale-cohomology/pinned-original-source.tex, source lines 2842--2842. Navigation pointer invoked at proof line 10733; the counted span contains the pointer only and the cited result is not reproduced inside it, so the pointer is registered as a non-load-bearing navigation stub in the established form for Stacks material. No bibliography citation is involved.\par\medskip
\ifdefined\hypertarget\hypertarget{sourcepointer-6}{}\fi
\noindent\textit{Source reference.} Original theorem 03P8: sources/stacks-etale-cohomology/pinned-original-source.tex, source lines 2498--2498. Navigation pointer invoked at proof line 10742; the counted span contains the pointer only and the cited result is not reproduced inside it, so the pointer is registered as a non-load-bearing navigation stub in the established form for Stacks material. No bibliography citation is involved.\par\medskip
\ifdefined\hypertarget\hypertarget{sourcepointer-7}{}\fi
\noindent\textit{Source reference.} Original lemma 0DV9: sources/stacks-etale-cohomology/pinned-original-source.tex, source lines 10658--10658. Navigation pointer invoked at proof line 10754; the counted span contains the pointer only and the cited result is not reproduced inside it, so the pointer is registered as a non-load-bearing navigation stub in the established form for Stacks material. No bibliography citation is involved.\par\medskip
\ifdefined\hypertarget\hypertarget{sourcepointer-8}{}\fi
\noindent\textit{Source reference.} Original section 0A3J: sources/stacks-etale-cohomology/pinned-original-source.tex, source lines 9600--9600. Navigation pointer invoked at proof line 10759; the counted span contains the pointer only and the cited result is not reproduced inside it, so the pointer is registered as a non-load-bearing navigation stub in the established form for Stacks material. No bibliography citation is involved.\par\medskip
\ifdefined\hypertarget\hypertarget{sourcepointer-9}{}\fi
\noindent\textit{Source reference.} Original remark 03P6: sources/stacks-etale-cohomology/pinned-original-source.tex, source lines 2473--2473. Navigation pointer invoked at proof line 10767; the counted span contains the pointer only and the cited result is not reproduced inside it, so the pointer is registered as a non-load-bearing navigation stub in the established form for Stacks material. No bibliography citation is involved.\par\medskip
\ifdefined\hypertarget\hypertarget{sourcepointer-10}{}\fi
\noindent\textit{Source reference.} Original lemma 01XB: sources/stacks-etale-cohomology/pinned-original-source.tex, source lines 10770--10770. Navigation pointer invoked at proof line 10770; the counted span contains the pointer only and the cited result is not reproduced inside it, so the pointer is registered as a non-load-bearing navigation stub in the established form for Stacks material. No bibliography citation is involved.\par\medskip
\ifdefined\hypertarget\hypertarget{sourcepointer-11}{}\fi
\noindent\textit{Source reference.} Original lemma 0DV3: sources/stacks-etale-cohomology/pinned-original-source.tex, source lines 8979--8979. Navigation pointer invoked at proof line 10790; the counted span contains the pointer only and the cited result is not reproduced inside it, so the pointer is registered as a non-load-bearing navigation stub in the established form for Stacks material. No bibliography citation is involved.\par\medskip
\begin{proposition}
\label{proposition-serre-galois}

Let $K$ be a field with separable algebraic closure $K^{sep}$.
Assume that for any finite extension $K'$ of $K$ we have
$\text{Br}(K') = 0$. Then
\begin{enumerate}
\item $H^q(\text{Gal}(K^{sep}/K), (K^{sep})^*) = 0$
for all $q \geq 1$, and
\item $H^q(\text{Gal}(K^{sep}/K), M) = 0$
for any torsion $\text{Gal}(K^{sep}/K)$-module $M$ and any $q \geq 2$,
\end{enumerate}
\end{proposition}

\begin{proof}
Set $p = \text{char}(K)$.
By Lemma \ref{lemma-compare-cohomology-point},
Theorem \ref{theorem-brauer-delta},
and Example \ref{example-sheaves-point}
the proposition is equivalent to showing that if 
$H^2(\Spec(K'),\mathbf{G}_m|_{\Spec(K')_\etale}) = 0$
for all finite extensions $K'/K$ then:
\begin{itemize}
\item $H^q(\Spec(K),\mathbf{G}_m|_{\Spec(K)_\etale}) = 0$
for all $q \geq 1$, and
\item $H^q(\Spec(K),\mathcal{F}) = 0$
for any torsion sheaf $\mathcal{F}$ and any $q \geq 2$.
\end{itemize}
We prove the second part first.
Since $\mathcal{F}$ is a torsion sheaf, we may use the 
$\ell$-primary decomposition as well as the compatibility of
cohomology with colimits (i.e, direct sums, see Theorem \ref{theorem-colimit})
to reduce to showing  $H^q(\Spec(K),\mathcal{F}) = 0$, $q \geq 2$
for all $\ell$-power torsion sheaves for every prime $\ell$. This 
allows us to analyze each prime individually.

\medskip\noindent
Suppose that $\ell \neq p$. For any extension $K'/K$
consider the Kummer sequence (Lemma \ref{lemma-kummer-sequence})
$$
0 \to
\mu_{\ell, \Spec{K'}} \to
\mathbf{G}_{m, \Spec{K'}} \xrightarrow{(\cdot)^{\ell}}
\mathbf{G}_{m, \Spec{K'}} \to 0
$$
Since $H^q(\Spec{K'},\mathbf{G}_m|_{\Spec(K')_\etale}) = 0$
for $q = 2$ by assumption  and for $q = 1$ by
Theorem \ref{theorem-picard-group} combined with
$\Pic(K) = (0)$. Thus, by the long-exact cohomology 
sequence we may conclude that $H^2(\Spec{K'}, \mu_\ell) = 0$ for 
any separable $K'/K$. Now let $H$ be a maximal pro-$\ell$ subgroup
of the absolute Galois group of $K$ and let $L$ be the 
corresponding extension. We can write $L$ as the colimit of finite extensions, 
applying Theorem \ref{theorem-colimit} to this colimit we see that
$H^2(\Spec(L), \mu_\ell) = 0$.
Now $\mu_\ell$ must be the constant sheaf. If it weren't, that would imply
there exists a Galois extension of degree relatively prime to 
$\ell$ of $L$ which is not true by definition of $L$ (namely, the extension
one gets by adjoining the  $\ell$th roots of unity to $L$).
Hence, via Lemma \ref{lemma-reduce-to-l-group-higher},
we conclude the result for $\ell \neq p$. 

\medskip\noindent
Now suppose that $\ell = p$. We consider the
Artin-Schrier exact sequence (Section \ref{section-artin-schreier})
$$
0 \longrightarrow \underline{\mathbf{Z}/p\mathbf{Z}}_{\Spec{K}} \longrightarrow
\mathbf{G}_{a, \Spec{K}} \xrightarrow{F-1} \mathbf{G}_{a, \Spec{K}} 
\longrightarrow 0
$$
where $F - 1$ is the map $x \mapsto x^p - x$. Then note that the higher 
Cohomology of $\mathbf{G}_{a, \Spec{K}}$ vanishes, by
Remark \ref{remark-special-case-fpqc-cohomology-quasi-coherent} and the 
vanishing of the higher cohomology of the structure sheaf of an affine scheme
(Cohomology of Schemes, Lemma
\ref{coherent-lemma-quasi-coherent-affine-cohomology-zero}).
Note this can be applied to any field of 
characteristic $p$. In particular, we can apply it to the field extension $L$ 
defined by a maximal pro-$p$ subgroup $H$. This allows us to conclude 
$H^n(\Spec{L}, \underline{\mathbf{Z}/p\mathbf{Z}}_{\Spec{L}}) = 0$
for $n \geq 2$, from which the result follows for $\ell = p$, by
Lemma \ref{lemma-reduce-to-l-group-higher}.

\medskip\noindent
To finish the proof we still have to show that
$H^q(\text{Gal}(K^{sep}/K), (K^{sep})^*) = 0$ for all $q \geq 1$.
Set $G = \text{Gal}(K^{sep}/K)$ and set $M = (K^{sep})^*$
viewed as a $G$-module. We have already shown (above) that
$H^1(G, M) = 0$ and $H^2(G, M) = 0$. Consider the exact sequence
$$
0 \to A \to M \to M \otimes \mathbf{Q} \to B \to 0
$$
of $G$-modules. By the above we have $H^i(G, A) = 0$
and $H^i(G, B) = 0$ for $i > 1$ since $A$ and $B$ are
torsion $G$-modules. By
Lemma \ref{lemma-profinite-group-cohomology-torsion}
we have $H^i(G, M \otimes \mathbf{Q}) = 0$ for $i > 0$.
It is a pleasant exercise to see that this implies that
$H^i(G, M) = 0$ also for $i \geq 3$.
\end{proof}

\end{document}
